\documentclass[11pt]{article}
\usepackage[margin=1.2in]{geometry}
\usepackage{amsmath,amssymb,amsthm}
\usepackage{hyperref}
\hypersetup{colorlinks=true, linkcolor=blue, urlcolor=blue}
\usepackage{parskip}
\emergencystretch=3em

\newcommand{\R}{\mathbb{R}}
\newcommand{\code}[1]{\texttt{#1}}

\title{The forward-expansion programme: a synthesis\\
(germbij Theorem 3.1's forward direction, twelve stages)}
\author{Tide loop (Claude Fable 5, with GPT-5.6 Sol deliberation)}
\date{2026-08-10}

\begin{document}
\maketitle

\begin{abstract}
Twelve stacked tides (PRs \#96--\#110) carried the germbij note's
Theorem~3.1 from its inverse half --- superpolynomially-close
localized moment families determine the location, the positive-order
jet, and for analytic losses the germ --- to the full nondegenerate
statement: on a \code{ForwardExpansionDomain} (a nondegenerate
minimum with $C^{N+2}$ regularity and a Peano remainder), the
rescaled posterior moment of every continuous observable of
polynomial growth \emph{is} an order-$N$ asymptotic polynomial at
$0^+$ (\code{rescaledMoment\_hasExpansion}), with monomial instances
matching the data consumed by the inverse half and a jet-comparison
corollary. Expansion data and Taylor data now determine each other in
both directions, in Lean. The programme followed a single up-front
design consult almost to the letter; its one architectural deviation
--- replacing the prescribed log-derivative coefficient recursion by
exp's Taylor truncation composed with the exponent polynomial ---
removed every coefficient identity from the analytic spine, and its
one analytic correction --- an unrestricted exponential remainder in
place of Mathlib's $|x| \le 1$ bound, with the $q$-power recovered
multinomially --- was forced by the geometry of the mesoscopic
window, where corrections are small only relative to the quadratic.
\end{abstract}

\section{What was proven}

\begin{itemize}
\item \textbf{AsymptoticPolynomial} (\#96): the expansion predicate
\code{IsAsymptoticExpansionTo} at $0^+$ and coefficient uniqueness ---
the tool that later closed both the quadratic bridge and the
jet-comparison corollary.
\item \textbf{GaussianMeso} (\#97): the mesoscopic window
$\sqrt q\,\|z\| \le 1$ and its tail calculus --- outside the window,
polynomial-Gaussian and rescaled-Boltzmann integrals are beyond all
orders.
\item \textbf{ForwardDomain} (\#98): the one-field Peano mixin over
the merged higher-order package, and the exact exponent split
$(L(qz)-L(0))/q^2 = T_2(z) + \sum_{s=1}^N q^s V_s(z) + q^N\rho_q(z)$
with pointwise vanishing and window bounds on $\rho$.
\item \textbf{ExpGraded} (\#99): the graded expansion of
$\exp(-(\sum_s q^s a_s + q^N \rho_q))$ with coefficients from the
Taylor-truncation polynomial --- three elementary steps, no
coefficient identities.
\item \textbf{CoeffFn} (\#100): continuity and polynomial growth of
the coefficient functions $z \mapsto P_j(z)$ at arbitrary degree, by
induction through the coefficient convolution.
\item \textbf{ScalarBounds} (\#101): the unrestricted exponential
Taylor remainder $|e^x - S_N(x)| \le |x|^{N+1}e^{|x|}/(N{+}1)!$, the
perturbation bound, and the graded polynomial tail with a
$z$-uniform indexing bound.
\item \textbf{GaussAbsorb} (\#102): $T_1 = 0$ and
$T_2 = \tfrac12\,\mathrm{qform}\,H$ as \emph{theorems} (ray
coefficient uniqueness), $\varepsilon$-Peano window control, and the
absorption of all corrections into $e^{-(\lambda/4)\|z\|^2}$.
\item \textbf{WindowMajorant} (\#103): the single $q$-uniform
majorant the dominated-convergence step consumes.
\item \textbf{NumeratorExpansion} (\#104) and \textbf{NumeratorTails}
(\#106): the window convergence by filter-indexed dominated
convergence with the indicator folded in; the two outer tails; and
\code{numerator\_hasExpansion} with coefficients
$\int P\, e^{-T_2} P_j$.
\item \textbf{AsymptoticDivision} (\#108): quotients of order-$N$
expansions with nonzero constant denominator coefficient.
\item \textbf{ForwardTheorems} (\#110): partition positivity, the
moment expansion, the monomial instances, and jet comparison.
\end{itemize}

\section{Design decisions that carried the programme}

\textbf{Existence-plus-comparison as the public API.} The programme
never needed explicit coefficient values: the note's theorem asserts
that expansions exist and determine the jet, and stage A1's
uniqueness converts any two expansion claims into coefficient
equality. Every collection mechanism stayed internal.

\textbf{The Taylor-truncation collection mechanism.} The design
consult prescribed the log-derivative recursion
$P_j = -\tfrac1j\sum s\,V_s P_{j-s}$ but noted the proof may use
exp's ordinary Taylor remainder. Defining the coefficients as the
coefficients of $\sum_{i\le N} (-A)^i/i!$ made the scalar expansion
theorem three elementary $O(q^{N+1})$ steps and moved all
combinatorics into a measurability-and-growth induction through
\code{Polynomial.coeff\_mul} --- the same mechanism then reused for
the division stage's Cauchy product.

\textbf{The unrestricted remainder over a second window.} On the
mesoscopic window the exponent correction reaches $\sim 1/\sqrt q$;
Mathlib's \code{Real.exp\_bound} ($|x| \le 1$) cannot see it. The
stage-5 consult confirmed the fix: prove
$|e^x - S_N(x)| \le |x|^{N+1}e^{|x|}/(N{+}1)!$ once, and recover the
$q$-power from $|A| \le q\sum_s|V_s|$ on $q \le 1$ rather than from
absolute smallness. A smaller window would have needed its own tail
theorems and still failed at $N = 0$.

\textbf{The quadratic bridge by self-consumption.} The domain ties
$H$ to $L$ only through the quadratic Peano condition, so
$T_2 = \tfrac12\,\mathrm{qform}\,H$ was proven by composing both
Peano statements along rays and applying the programme's own stage-1
uniqueness --- which also yielded the critical-point condition
$T_1 = 0$, converting the exponent split's hypothesis into a theorem.

\textbf{One majorant, one DCT.} All domain-specific uniformity
concentrated in a single theorem
(\code{normalized\_window\_remainder\_bound}); the dominated
convergence then ran with the window indicator folded into the
integrand, per the consult, and the final packaging was four
subidentities closed by linear arithmetic.

\section{Process notes}

Two architecture consults (the seven-stage programme design; the
stage-5 shape with its nine-step implementation order) were followed
essentially verbatim, and both are archived in the tide log. Every
stage's interface was consumed by the next without rework; the
majorant --- the longest single proof --- compiled on its second
build attempt. Recurring Lean friction was almost entirely
catalogued classes: beta-redex goals from \code{integral\_congr\_ae}
(now: \code{beta\_reduce} first), norms on both sides of
\code{isLittleO\_iff}, \code{positivity} blind to opaque structure
fields, and two deterministic whnf heartbeat timeouts. New catalogue
entries from this arc are recorded in the per-tide retrospectives.

\section{What remains}

Of the germbij note: the semi-quasi-homogeneous constructive class
(the $M_\alpha$ recovery, Gelfand--Leray, and the weighted triangular
induction), ranked less crucial in the commissioning direction; and,
if ever wanted, the connection of the abstract \code{momentCoeff} to
explicit Wick/Isserlis values through the Gaussian moment integrals
already in the seabed. A parallel fidelity review (PR \#105) flags
the inverse-half headline as conditional on a matched $0$--$2$ jet;
its perimeter debt list is the right starting point for any
tightening pass.
\end{document}
